\documentclass{article}
\usepackage{mathtools} 
\usepackage{fontspec}
\usepackage[UTF8]{ctex}
\usepackage{amsthm}
\usepackage{mdframed}
\usepackage{xcolor}
\usepackage{amssymb}
\usepackage{amsmath}


% 定义新的带灰色背景的说明环境 zremark
\newmdtheoremenv[
  backgroundcolor=gray!10,
  % 边框与背景一致，边框线会消失
  linecolor=gray!10
]{zremark}{说明}

% 通用矩阵命令: \flexmatrix{矩阵名}{元素符号}{行数}{列数}
\newcommand{\flexmatrix}[4]{
  \[
  #1 = \begin{pmatrix}
    #2_{11}     & #2_{12}     & \cdots & #2_{1#4}   \\
    #2_{21}     & #2_{22}     & \cdots & #2_{2#4}   \\
    \vdots      & \vdots      & \ddots & \vdots     \\
    #2_{#31}    & #2_{#32}    & \cdots & #2_{#3#4}
  \end{pmatrix}
  \]
}

% 简化版命令（默认矩阵名为A，元素符号为a）: \quickmatrix{行数}{列数}
\newcommand{\quickmatrix}[2]{\flexmatrix{A}{a}{#1}{#2}}

\begin{document}
\title{习题 3.3}
\author{张志聪}
\maketitle

\section*{1}

设$f$在$U(+\infty)$上有定义。$\lim\limits_{x \to +\infty}f(x)$
存在的充要条件是：对任何含于$U(+\infty)$且以$+\infty$为极限的数列
$\{x_n\}$，极限$\lim\limits_{n \to +\infty}f(x_n)$都存在且相等。

\section*{2}

只证明增函数的情况，减函数的情况同理。

\begin{itemize}
  \item 必要性

        不妨设$\lim\limits_{x \to +\infty}f(x) = A$，
        于是$\epsilon = 1, \exists M$当$x > M$时，有
        \begin{align*}
          |f(x) - A| < \epsilon \\
          A - \epsilon < f(x) < A + \epsilon
        \end{align*}
        由$f$在$[a, +\infty)$上单调递增，所以
        $\forall x \in [a, M]$，我们有
        \begin{align*}
          f(a) \leq f(x) \leq f(M)
        \end{align*}
        综上，对$\forall x \in [a, +\infty)$，我们有
        \begin{align*}
          f(a) \leq f(x) \leq A + \epsilon
        \end{align*}
        即$f$在$[a, +\infty)$上有上界。

  \item 充分性

        因为$f$在$[a, +\infty)$中有上界，
        由确界定理可知，$f$在$[a, +\infty)$有上确界。

        设$A = \sup\limits_{x \in [a, +\infty)} f(x)$，
        我们证明$\lim\limits_{x \to +\infty} f(x) = A$。

        由上确界的定义可知，$\forall \epsilon > 0$，存在$x^\prime \in [a, +\infty)$，
        使得
        \begin{align*}
          A - \epsilon < f(x^\prime) < A
        \end{align*}
        取$M = x^\prime$，则由$f$的递增性，对$x \geq M$，有
        \begin{align*}
          A - \epsilon < f(x^\prime) \leq f(x) \leq A < A + \epsilon
        \end{align*}
        所以$\lim\limits_{x \to +\infty} f(x) = A$。
\end{itemize}

\section*{4}

反证法，假设存在数列
\begin{align*}
  \lim\limits_{n \to +\infty} a_n = x_0 \\
  \lim\limits_{n \to +\infty} b_n = x_0
\end{align*}
且
\begin{align*}
  \lim\limits_{n \to +\infty} f(a_n) \neq \lim\limits_{n \to +\infty} f(b_n)
\end{align*}

构造数列
\begin{align*}
  \{c_n\} = \{a_1, b_1, a_2, b_2, \cdots, a_n, b_n, \cdots\}
\end{align*}
易得
\begin{align*}
  \lim\limits_{n \to +\infty} c_n = x_0
\end{align*}
有题设可知
\begin{align*}
  \lim\limits_{n \to +\infty} f(c_n)
\end{align*}
是存在的，但数列$\{f(c_n)\}$的子列$\{f(a_n)\}, \{f(a_n)\}$
的极限不相等，可得$\lim\limits_{n \to +\infty} f(c_n)$不存在，
于是产生矛盾，假设不成立，命题得证。

\section*{6}

利用归结原理证明。

取一个有理数数列$\{x_n\}$，收敛于$x_0$，即
\begin{align*}
  \lim\limits_{n \to +\infty} x_n = x_0
\end{align*}

取一个无理数数列$\{y_n\}$，收敛于$y_0$，即
\begin{align*}
  \lim\limits_{n \to +\infty} y_n = x_0
\end{align*}

于是，我们有
\begin{align*}
  \lim\limits_{n \to +\infty} f(x_n) = 1 \\
  \lim\limits_{n \to +\infty} f(y_n) = 0
\end{align*}

所以$\lim\limits_{x \to x_0} D(x)$不存在。

这里需要补充一下数列$\{x_n\}, \{y_n\}$的存在性。\\  
在$U^{\circ}(x_0; 1)$上，根据实数的稠密性，可取有理数$x_1$和无理数$y_1$，
此时
\begin{align*}
  x_0 - 1 < x_1 < x_0 + 1 \\
  x_0 - 1 < y_1 < x_0 + 1
\end{align*}
在$U^{\circ}(x_0; \frac{1}{2})$上，根据实数的稠密性，可取有理数$x_2$和无理数$y_2$，
此时
\begin{align*}
  x_0 - \frac{1}{2} < x_2 < x_0 + \frac{1}{2} \\
  x_0 - \frac{1}{2} < y_2 < x_0 + \frac{1}{2} 
\end{align*}
$\vdots$\\
在$U^{\circ}(x_0; \frac{1}{n})$上，根据实数的稠密性，可取有理数$x_n$和无理数$y_n$，
此时
\begin{align*}
  x_0 - \frac{1}{n} < x_n < x_0 + \frac{1}{n} \\
  x_0 - \frac{1}{n} < y_n < x_0 + \frac{1}{n} 
\end{align*}
$\vdots$\\
由迫敛性可知
\begin{align*}
  \lim_{n\to\infty} x_n = x_0 \quad \lim_{n\to\infty} y_n = x_0
\end{align*}

\end{document}